\documentclass[11pt,a4paper]{article}

% ---------------------------------------------------------
% PREAMBLE (stable, warning-free, Overleaf-safe)
% ---------------------------------------------------------
\usepackage[utf8]{inputenc}
\usepackage[T1]{fontenc}
\usepackage{lmodern}
\usepackage{amsmath, amssymb, amsfonts}
\usepackage{bm}
\usepackage{geometry}
\usepackage{graphicx}
\usepackage{hyperref}
\usepackage{cite}
\usepackage{authblk}
\usepackage{physics}
\usepackage{mathtools}

\geometry{margin=2.4cm}

\title{\textbf{Companion LXXVI: Wasserstein Information Geometry of Persistence Diagrams in the Relaxation-Driven Cyclic Cosmology}}
\author{Michael Lehmann \\ \texttt{mi.lehmann@gmx.de}}
\date{April 2026}

\begin{document}
\maketitle

\begin{abstract}
Wasserstein geometry provides the natural metric structure on the space of 
persistence diagrams, capturing the optimal transport cost required to match 
topological features across diagrams. Information geometry, in contrast, equips 
statistical manifolds with a Riemannian structure derived from Fisher information. 
This Companion develops a unified Wasserstein information geometry for persistence 
diagrams in the Relaxation-Driven Cyclic Cosmology (RDCC). The relaxation 
mechanism modifies the scale-dependent evolution of the gravitational potential, 
altering both the Wasserstein transport structure and the information-theoretic 
curvature of topological ensembles. The resulting framework provides a 
geometrically rich and observationally relevant description of RDCC signatures in 
weak-lensing topology.
\end{abstract}
% ---------------------------------------------------------
\section{Introduction}

Persistence diagrams encode the birth and death of topological features in weak-
lensing convergence fields. Wasserstein distances provide the canonical metric on 
the space of diagrams, measuring the optimal transport cost required to match 
topological features across realizations. Previous Companions developed Hilbert-
space embeddings, kernel PCA, MMD, and Fisher information geometry as tools for 
analyzing the Relaxation-Driven Cyclic Cosmology (RDCC).

Wasserstein information geometry synthesizes these perspectives by equipping the 
Wasserstein space of persistence diagrams with an information-theoretic Riemannian 
structure. This structure quantifies how sensitively the optimal transport 
geometry responds to changes in cosmological parameters. In RDCC, the relaxation 
mechanism modifies the gravitational potential in a scale-dependent manner, 
producing characteristic signatures in the Wasserstein information geometry of 
topological ensembles.
% ---------------------------------------------------------
\section{Wasserstein Geometry of Persistence Diagrams}

Given two persistence diagrams $D$ and $E$, the $p$-Wasserstein distance is

\begin{equation}
    W_{p}(D,E)
    = \left(
        \inf_{\gamma}
        \sum_{x \in D}
        \| x - \gamma(x) \|^{p}
      \right)^{1/p},
\end{equation}

where $\gamma$ ranges over bijections between $D$ and $E$ (including the diagonal).

This metric captures the minimal transport cost required to match topological 
features across diagrams.
% ---------------------------------------------------------
\section{Information Geometry on Wasserstein Space}

Let $p(D|\theta)$ be a statistical model for persistence diagrams. The Fisher 
information matrix is

\begin{equation}
    I_{ab}(\theta)
    = \mathbb{E}
      \left[
        \partial_{\theta_{a}} \log p(D|\theta)
        \,
        \partial_{\theta_{b}} \log p(D|\theta)
      \right].
\end{equation}

The Wasserstein information metric is defined by pulling back the Fisher metric 
through the Wasserstein map:

\begin{equation}
    g^{W}_{ab}(\theta)
    = \mathbb{E}
      \left[
        \partial_{\theta_{a}} W_{p}(D, \mu(\theta))
        \,
        \partial_{\theta_{b}} W_{p}(D, \mu(\theta))
      \right],
\end{equation}

where $\mu(\theta)$ is the Wasserstein barycenter of the distribution.
% ---------------------------------------------------------
\section{RDCC Modification of Wasserstein Structure}

In RDCC, the convergence evolves as

\begin{equation}
    \kappa_{\mathrm{RDCC}}(\ell)
    = \kappa(\ell)
      \left[
        1 - \chi_{\kappa}
        \left(\frac{\ell}{\ell_{\mathrm{rel}}}\right)^{\alpha}
      \right].
\end{equation}

This modifies the birth and death coordinates of topological features:

\begin{align}
    b_{\mathrm{RDCC}} &= b \left[ 1 - \chi_{b} (\ell/\ell_{\mathrm{rel}})^{\alpha} \right], \\
    d_{\mathrm{RDCC}} &= d \left[ 1 - \chi_{d} (\ell/\ell_{\mathrm{rel}})^{\alpha} \right].
\end{align}

The Wasserstein distance becomes

\begin{equation}
    W_{p}^{\mathrm{RDCC}}(D,E)
    = W_{p}(D,E)
      \left[
        1 - \chi_{W}
        \left(\frac{\ell}{\ell_{\mathrm{rel}}}\right)^{\alpha}
      \right].
\end{equation}
% ---------------------------------------------------------
\section{Wasserstein Information Geometry in RDCC}

The RDCC Wasserstein information metric is

\begin{equation}
    g^{W,\mathrm{RDCC}}_{ab}
    = g^{W}_{ab}
      \left[
        1 - \chi_{g,ab}
        \left(\frac{\ell}{\ell_{\mathrm{rel}}}\right)^{\alpha}
      \right].
\end{equation}

RDCC induces:

\begin{itemize}
    \item reduced Wasserstein curvature at small scales,
    \item enhanced curvature at large scales,
    \item a characteristic turnover near $\ell_{\mathrm{rel}}$,
    \item anisotropic deformation of Wasserstein geodesics,
    \item modified barycentric flow of topological ensembles.
\end{itemize}
% ---------------------------------------------------------
\section{Conclusion}

Wasserstein information geometry provides a powerful framework for analyzing the 
optimal-transport sensitivity of persistence diagrams in the Relaxation-Driven 
Cyclic Cosmology. The relaxation mechanism modifies both the Wasserstein transport 
structure and the information-theoretic curvature of topological ensembles in a 
scale-dependent manner. These effects offer a geometrically rich and 
observationally relevant probe of the RDCC gravitational field. The present 
Companion establishes the foundation for future studies of Wasserstein geometry, 
information geometry, and optimal transport in cosmological topology.

\bibliographystyle{apsrev4-2}
\nocite{*}
\bibliography{references}


\end{document}
